\begin{frame}
        \frametitle{Angular Discretization}
        Before applying any discretizations, we first examine the streaming term (\cref{eq:streaming-term}).
        By aligning the polar angle with the $x$-axis, the solid angle is written in Cartesian coordinates:
        \begin{equation}
                \label{eq:solid-angle-expansion}
                \ohat = \underbrace{\cos \theta}_{\ohat_x} \mathbf{\hat{x}}+\underbrace{\sin\theta \sin\varphi}_{\ohat_y} \mathbf{\hat{y}} + \underbrace{\sin\theta \cos \varphi}_{\ohat_z} \mathbf{\hat{z}}
        \end{equation}
        Substituting into \cref{eq:streaming-term}:
        \begin{equation}
                \begin{split}
                        \ohat \cdot \grad \psisimple & = \ohat_x \pdv{\psisimple}{x}+\ohat_y\pdv{\psisimple}{y}+\ohat_z\pdv{\psisimple}{z} \\
                                                     & = \mu \pdv{\psisimple}{x}+\eta\pdv{\psisimple}{y}+\xi\pdv{\psisimple}{z}
                \end{split}
        \end{equation}
        where, in the second line, we have substituted in greek letters to fit with convention.
\end{frame}

\begin{frame}[allowframebreaks]
        \frametitle{Angular Quadrature}
        We look to approximate the integral over the unit sphere for the angle-integrated flux:
        \begin{equation}\label{eq:scalar-flux-integral-expanded}
                \begin{split}
                        \int_{4\pi}\psisimple\dd \ohat & = \int_{0}^{2\pi}\int_{0}^{\pi}\psi\left(\rvec,\theta,\varphi\right)\dd \theta \dd \varphi \\
                                                       & =\int_{0}^{2\pi}\int_{-1}^{1}\psi\left(\rvec,\mu,\varphi\right)\dd \mu \dd \varphi         \\
                                                       & = \int_{0}^{2\pi}\Gamma\left(\rvec,\varphi\right) \dd \varphi
                \end{split}
        \end{equation}
        where $\mu=\cos\theta$ as before and $\Gamma\left(\rvec,\varphi\right)\equiv\int_{-1}^{1}\psi\left(\rvec,\mu,\varphi\right)\dd \mu$.
        \Cref{eq:scalar-flux-integral-expanded} can be approximated with two quadrature sets:
        \begin{subequations}
                \begin{equation}\label{eq:polar-quad}
                        \Gamma\left(\rvec,\varphi\right)\equiv\int_{-1}^{1}\psi\left(\rvec,\mu,\varphi\right)\dd \mu \approx \sum_{m=1}^M w_m \psi\left(\rvec,\mu_m,\varphi\right)
                \end{equation}
                \begin{equation}\label{eq:azim-quad}
                        \int_{0}^{2\pi}\Gamma\left(\rvec,\varphi\right) \dd \varphi \approx \sum_{n=1}^N w_n \Gamma\left(\rvec,\varphi_n\right)
                \end{equation}
        \end{subequations}
        Combining \cref{eq:polar-quad} and \cref{eq:azim-quad}:
        \begin{equation}
                \label{eq:composite-quad}
                \int_{4\pi}\psisimple\dd \ohat \approx \sum_{n=1}^N \sum_{m=1}^M w_{nm}\psi\left(\rvec,\mu_m,\varphi_n\right)
        \end{equation}

        In this implementation, \cref{eq:polar-quad} is approximated with Gauss-Legendre quadrature, while \cref{eq:azim-quad} is approximated with Gauss-Chebyshev quadrature.
        The discretization of angle generates a set of \textit{discrete ordinates}, thus we look to solve:
        \begin{equation}\label{eq:sn-equations}
                \mu_{m}\dv{\psi_{mn}}{x}+\eta_{n}\dv{\psi_{mn}}{y} + \Sigma_t \left(\rvec\right)\psi_{mn}\left(\rvec\right) =\frac{\Sigma_s}{4\pi}\phi\left(\rvec\right) + \frac{Q\left(\rvec\right)}{4\pi}
        \end{equation}
        where $\psi_{mn}\left(\rvec\right)\equiv \psi\left(\rvec,\mu_m,\varphi_n\right)$, $\mu=\cos\theta$, and $\eta=\sqrt{1-\mu^2}\sin\varphi$.
        \Cref{eq:sn-equations} is a set of $N\times M$ equations coupled by \cref{eq:composite-quad}.
\end{frame}
